\documentclass[11pt]{amsart}
\usepackage{calc,amssymb,amsmath,ulem,stmaryrd,fullpage}
\usepackage{enumerate}
\usepackage{alltt}
\RequirePackage[dvipsnames,usenames]{color}
\let\mffont=\sf
\normalem
%\input{kmacros3.sty}
\input{xy}
\xyoption{all}
\usepackage{hyperref}
\usepackage{graphicx}
\usepackage[all,cmtip]{xy}
\usepackage{verbatim}
\usepackage[left=0.95in,top=0.95in,right=0.95in,bottom=0.95in]{geometry}
\newcommand{\tauCohomology}{T}
\newcommand{\FCohomology}{S}


\theoremstyle{theorem}
%\newtheorem*{mainthm*}{Main Theorem}

\newcommand{\note}[1]{\marginpar{\sffamily\tiny #1}}

\def\todo#1{\textcolor{Mahogany}%
{\sffamily\footnotesize\newline{\color{Mahogany}\fbox{\parbox{\textwidth-15pt}{#1}}}\newline}}

\renewcommand{\O}{\mathcal O}
\newcommand{\ie}{{\itshape i.e.} }
\newcommand{\roundup}[1]{\lceil #1 \rceil}
\newcommand{\rounddown}[1]{\lfloor #1 \rfloor}
\renewcommand{\to}[1][]{\xrightarrow{\ #1\ }}
\newcommand{\onto}[1][]{\protect{\xrightarrow{\ #1\ }\hspace{-0.8em}\rightarrow}}
\newcommand{\into}[1][]{\lhook \joinrel \xrightarrow{\ #1\ }}
%\newcommand{DBDual}[1]{{\underline{\omega}_{#1}}}
\newcommand{\DBDual}[1]{{{\uuline \omega}{}^{\mydot}_{#1}}}
\newcommand{\Tt}{{\mathfrak{T}}}
%\newcommand{\bn}{{\mathfrak{n}} }
\newcommand{\sol}[1]{ \vskip 6pt{\color{blue} {\bf Solution:} #1} \vskip 6pt}

\newcommand{\bR}{{\mathbb{R}}}
\newcommand{\va}{{\vec{a}}}
\newcommand{\vb}{{\vec{b}}}
\newcommand{\vzero}{{\vec{0}}}
\newcommand{\vi}{{\vec{i}}}
\newcommand{\vj}{{\vec{j}}}
\newcommand{\vk}{{\vec{k}}}
\newcommand{\vh}{{\vec{h}}}
\newcommand{\vr}{{\vec{r}}}
\newcommand{\vu}{{\vec{u}}}
\newcommand{\vv}{{\vec{v}}}
\newcommand{\vw}{{\vec{w}}}
\newcommand{\vx}{{\vec{x}}}
\newcommand{\vy}{{\vec{y}}}
\newcommand{\vz}{{\vec{z}}}
\newcommand{\vT}{{\vec{T}}}
\newcommand{\vN}{{\vec{N}}}
\newcommand{\vF}{{\vec{F}}}
\newcommand{\vG}{{\vec{G}}}
\newcommand{\bF}{\mathbb{F}}
\newcommand{\bQ}{\mathbb{Q}}
\newcommand{\bZ}{\mathbb{Z}}
\newcommand{\Inn}{\mathrm{Inn}}
\newcommand{\Aut}{\mathrm{Aut}}
\newcommand{\Ann}{\mathrm{Ann}}
\newcommand{\tor}{\mathrm{tor}}


\begin{document}
\title{Worksheet \#4 -- Math 6310\\Fall 2019}
\author{Due: Friday, October 18th}
\maketitle

\noindent
{\bf Definition: }  Suppose $R$ is a ring.  A \emph{left $R$-module} $M$ is an Abelian group under $+$ with a multiplication map by elements of $R$:
\[
    \xymatrix@R=1pt{
        R \times M \ar[r] & M\\
        (r, m) \ar@{|->}[r] & r.m
    }
\]
satisfying the following axioms.
\begin{itemize}
\item[(1)] $r_1.(r_2.m) = (r_1 r_2).m$ for all $r_i \in R$, $m \in M$.
\item[(2)] $1.m = m$ for all $m \in M$.
\item[(3)] $r(m+n) = rm + rn$ for all $r \in R$, $m,n \in M$.
\item[(4)] $(r_1 + r_2)m = r_1 m + r_2 m$ for all $r_i \in R$, $m \in M$.
\end{itemize}
We can define a right $R$-module likewise (there is no real difference if $R$ is commutative).
\vskip 9pt
It is important to note that if $R$ is a field, an $R$-module is nothing more than an $R$-vector space.

Another important example, the left ideals of $R$ are exactly the $R$-submodules of $R$ acting on itself.
\vskip 9pt
\noindent
{\bf 1.}  Suppose that $\phi : R \to S$ is a ring homomorphism and $M$ is any $S$-module.  Show that we can view $M$ as an $R$-module via $r.m = \phi(r).m$.
\newpage
\noindent
{\bf 2.}  Show that every Abelian group $G$ under $+$ is a $\bZ$-module.
\newpage
\noindent
{\bf 3.}  Suppose $R$ is a ring and $M$ is an $R$-module.  Define
\[
    \Ann_R(M) = \{r \in R\;|\; rm = 0 \;\forall m \in M \}.
\]
Prove that this is an ideal of $R$ and that we can also view $M$ as an $R/\Ann_R(M)$-module.
\newpage
\noindent
{\bf 4.}  Suppose that $N \subseteq M$ is a submodule of a module.  Define the $R$-module $M/N$ and prove the $R$-module action you define is well defined.
\newpage
\noindent
A homomorphism of (left) $R$-modules is a function $\phi : M \to N$ between $R$-modules that is a homomorphism of Abelian groups (under addition) and such that $\phi(r.m) = r.\phi(m)$.
\vskip 9pt
\noindent
{\bf 5.}  Formulate and prove the first isomorphism theorem for homomorphisms of left $R$-modules.
\newpage
\noindent
{\bf 6.}  Suppose that $A \subseteq B \subseteq C$ are left $R$-modules.  Prove that $C/B \cong (C/A)/(B/A)$ as left $R$-modules.
\newpage
\noindent
{\bf 7.}  Suppose $R$ is a commutative integral domain and that $M$ is an $R$-module.  Consider the set
\[
    M_{\tor} = \{m \in M \;|\; rm = 0 \text{ for some nonzero } r\in R\}.
\]
\begin{itemize}
    \item[(a)]  Prove that $M_{\tor}$ is a submodule of $M$ (called the \emph{torsion submodule} of $M$).
    \item[(b)]  If $M_{\tor} = \{0\}$ then $M$ is called \emph{torsion free}.  Prove that $M/M_{\tor}$ is torsion free.
\end{itemize}
\newpage
\noindent
{\bf 8.}  Suppose $R$ is a commutative integral domain, $W \subseteq R$ is a multiplicative set and $M$ is an $R$-module.
Consider the set of pairs $(m, w)$ with $m \in M$ and $w \in W$ and declare two pairs $(m, w) \sim (m', w')$ to be equivalent if there exists $v \in W$ such that $v w' m = v w m'$.
\begin{itemize}
    \item[(a)]  Show that this defines an equivalence relation.  The equivalence class of $(m,w)$ is denoted by $m/w$.
    \item[(b)]  Prove that $W^{-1} M$ is a $W^{-1}R$-module (and also an $R$-module via the map $R \to W^{-1}R$).
    \item[(c)]  If $R = \bZ$, $W = \{1, 2, 2^2, 2^3, \dots\}$ and $M = \bZ/6\bZ$.  Prove that $W^{-1} M$ is isomorphic to $\bZ/3\bZ$.
\end{itemize}

\end{document}

